\documentclass[../main.tex]{subfiles}

\begin{document}
% -----------------------------
\section{Analysis}
% -----------------------------
Purpose of analysis is to compute information about the IR without changing it,
that can drive later optimizations/codegen.
We will implement two classic analyses on a function body block:
\begin{itemize}
    \item Use-def report - for each SSA result, list which ops use it
\end{itemize}
\begin{itemize}
    \item Block-local liveness - compute the live set after each op in a straight-line block
\end{itemize}

For typing annotation, import additional packages:
\begin{lstlisting} [language=bash, caption={Import typing}]
from typing import Any, Iterable
\end{lstlisting}
Create a couple of helper function:
\begin{lstlisting} [language=bash, caption={Analysis helper functions}]
# ----------------------------------
# Analysis
# ----------------------------------

# Helper functions
def _iter_uses(val: SSAValue) -> Iterable[Any]:
    uses = getattr(val, "uses", None)
    if uses is None:
        return []
    return uses

def _use_owner_op(use: Any) -> Operation | None:
    """
    Try to get the user operation from a Use object.
    """
    for attr in ("operation", "op", "owner"):
        u = getattr(use, attr, None)
        if isinstance(u, Operation):
            return u
    # some Use objects store the operation under .operation on a field
    return None

def _fmt_val(v: SSAValue) -> str:
    try:
        return str(v)  # often prints as %0
    except Exception:
        return f"<ssa@{id(v):x}>"

def _fmt_op(op: Operation) -> str:
    # compact one-liner label
    return f"{op.name}"

def _val_key(v: SSAValue) -> int:
    """
    Use id() as a stable key even if SSAValue is not hashable in some builds.
    """
    return id(v)

def _collect_operands(op: Operation) -> list[SSAValue]:
    return list(getattr(op, "operands", []))

def _collect_results(op: Operation) -> list[SSAValue]:
    return list(getattr(op, "results", []))
\end{lstlisting}
\begin{itemize}
    \item \path{_iter_uses} — Every \texttt{SSAValue} knows which operations use it,
    and that information is stored in \texttt{val.uses}. This function returns an
    iterable, even if \texttt{.uses} does not exist.

    \item \path{_use_owner_op} — Extracts the operation that uses a given
    \texttt{SSAValue} in a defensive way that tolerates different \texttt{Use}
    representations.

    \item \path{_fmt_val} — Creates a safe, human-readable string representation
    of an \texttt{SSAValue}.

    \item \path{_fmt_op} — Provides a compact textual label for an operation.

    \item \path{_val_key} — Creates a stable, hashable identity key for an
    \texttt{SSAValue}, even if the \texttt{SSAValue} itself cannot be used as a
    dictionary key.

    \item \path{_collect_operands} — Collects an operation's operands as a plain
    Python list.

    \item \path{_collect_results} — Collects an operation's results as a concrete
    Python list.
\end{itemize}



Create a function for printing use-def report:
\begin{lstlisting} [language=bash, caption={Creating use-def report}]
# ---------------------------------
# Analysis 1: Use-def report
# ---------------------------------
def print_use_def(block: Block) -> None:
    print("\n--- Use-def report (per SSA result) ---")
    for op in block.ops:
        for res in _collect_results(op):
            users: list[str] = []
            for use in _iter_uses(res):
                uop = _use_owner_op(use)
                if uop is None:
                    users.append("<unknown-user>")
                else:
                    users.append(_fmt_op(uop))
            users_s = ", ".join(users) if users else "(no uses)"
            print(f"{_fmt_val(res)} defined by {_fmt_op(op)}  -> used by: {users_s}")
\end{lstlisting}

Create a function for printing liveness report:
\begin{lstlisting} [language=bash, caption={Creating block-local liveness report}]
# -------------------------------------------------
# Analysis 2: Block-local liveness (straight-line)
# -------------------------------------------------
def compute_liveness(block: Block):
    """
    Returns a list `live_after[i]`: dict of {id(ssa): ssa} live AFTER block.ops[i].
    """
    ops = list(block.ops)
    live_after: list[dict[int, SSAValue]] = [dict() for _ in ops]

    live: dict[int, SSAValue] = {}

    # Seed with values used by the terminator (usually func.return)
    if ops:
        term = ops[-1]
        for v in _collect_operands(term):
            live[_val_key(v)] = v

    # Walk backwards
    for i in range(len(ops) - 1, -1, -1):
        op = ops[i]

        # live_after for this op = current live (after processing later ops)
        live_after[i] = dict(live)

        # Remove defs (results) from live (killing)
        for res in _collect_results(op):
            live.pop(_val_key(res), None)

        # Add uses (operands) to live
        for v in _collect_operands(op):
            live[_val_key(v)] = v

    return live_after
\end{lstlisting}

Implement a function that will print liveness based on previously added \path{compute_liveness}.
\begin{lstlisting} [language=bash, caption={Creating a function for printing liveness}]
def print_liveness(block: Block) -> None:
    live_after = compute_liveness(block)

    print("\n--- Block liveness (live set AFTER each op) ---")
    ops = list(block.ops)
    for i, op in enumerate(ops):
        live_vals = sorted((_fmt_val(v) for v in live_after[i].values()))
        live_s = ", ".join(live_vals) if live_vals else "(empty)"
        print(f"[{i:02d}] after {_fmt_op(op):>12} : {live_s}")
\end{lstlisting}

Implement a function for printing report:
\begin{lstlisting} [language=bash, caption={Function body block analysis}]
def analyze(module: ModuleOp):
    # Analyze each function body block
    for top_block in module.body.blocks:
        for op in top_block.ops:
            if isinstance(op, func.FuncOp):
                print(f"\n============================")
                print(f"Analyzing Function: {op.sym_name.data if hasattr(op, 'sym_name') else 'func'}")
                print(f"============================")

                for body_block in op.body.blocks:
                    print_use_def(body_block)
                    print_liveness(body_block)
\end{lstlisting}
and call it in the main function, after lowering and optimizations:
\begin{lstlisting} [language=bash, caption={\texttt{analyze} call in the main}]
...
analyze(module)
\end{lstlisting}
Disable constant folding and dead code elimination and run script again. Observe the results.

\end{document}
